% Short abstract ("shortest version for NBT" in the Google Doc); not included by main.tex
\section*{Abstract}
Long-read transcriptome sequencing enables isoform profiling in bulk, single-cell and single-nucleus samples, but partial transcripts and ambiguous read assignments impede accurate isoform discovery and quantification, particularly when transcript annotations are incomplete. We present the Long Read Alignment Assembler (LRAA), which integrates splice-graph-based structural modeling with expectation-maximization inference to discover isoforms and estimate their abundance. LRAA supports quantification-only, reference-guided and annotation-free modes within a unified framework. Across simulated and experimental PacBio and Oxford Nanopore benchmarks, including an innovative benchmark based on 3k biologically expressed, barcoded isoforms with read-level ground truth, LRAA ranked among the top-performing methods for isoform identification and quantification. Applied to single-cell and single-nucleus data, a cluster-guided workflow recovered cell-type-restricted isoforms in peripheral blood immune cells and detected a disease-associated cryptic STMN2 isoform in frontal cortex tissue from an individual with frontotemporal dementia. Transcript-derived open reading frame predictions suggested that approximately 13\% of predicted PBMC proteoforms and 5\% of those from frontal cortex arose exclusively from unannotated splice junctions or intron retention.
